% Generated draft: requires mathematical and source review.
\documentclass[11pt]{article}
\usepackage[margin=25mm]{geometry}
\usepackage{fontspec}
\setmainfont{Latin Modern Roman}
\usepackage{amsmath,amssymb,mathtools,unicode-math}
\setmathfont{Latin Modern Math}
\usepackage{enumitem,needspace,xurl,hyperref}
\setlength{\parindent}{0pt}
\setlength{\parskip}{5pt}
\setlength{\emergencystretch}{3em}
\newcommand{\N}{\mathbb{N}}
\newcommand{\Z}{\mathbb{Z}}
\newcommand{\Q}{\mathbb{Q}}
\newcommand{\R}{\mathbb{R}}
\newcommand{\C}{\mathbb{C}}
\newcommand{\F}{\mathbb{F}}
\newcommand{\A}{\mathbb{A}}
\newcommand{\PP}{\mathbb{P}}
\newcommand{\E}{\mathbb{E}}
\newcommand{\eps}{\varepsilon}
\newcommand{\dd}{\,\mathrm{d}}
\begin{document}
\section*{UnsolvedMath ID: 30002072}
\textbf{Source problem code:} catalogue entry
\subsection{Definitions and mathematical statement}
The supplied problem concerns the Cheeger--Müller theorem for singular manifolds and, in the accompanying literature assessment, a singular-space gluing theorem. Let \(Y\) denote the gluing boundary. A product-near-boundary hypothesis requires a collar identification with \(Y\times[0,\varepsilon)\), for some \(\varepsilon>0\), under which the structures covered by that hypothesis have product form. The supplied record does not specify the class of singular spaces, the structures concerned, the analytic domains or boundary conditions, or the precise theorem and gluing identity.

\textbf{Question (removal of the product hypothesis).} Does the cited singular-space Cheeger--Müller gluing result remain valid when the requirement that all relevant structures be products near the gluing boundary is omitted, with every other hypothesis retained?

More formally, let \(D\) stand for the complete data of the intended theorem, let \(H_{\mathrm{rest}}(D)\) denote all its hypotheses other than the product-near-boundary requirement, and let \(C(D)\) denote its conclusion. The requested extension has the schematic form
\[
\forall D,\qquad H_{\mathrm{rest}}(D)\Longrightarrow C(D).
\]
Here the range of \(D\), the hypotheses \(H_{\mathrm{rest}}\), and the conclusion \(C\) must be taken from the intended source theorem; they are placeholders, not additional mathematical assumptions. The supplied record alone does not determine them sufficiently to give a fully explicit quantified theorem. In particular, it does not resolve whether removing the product hypothesis is intended to preserve the identical gluing formula or to allow additional correction terms.
\subsection{Short English statement}
Can the product-near-boundary assumption be removed from the singular-space gluing result associated with the Cheeger--Müller theorem, while retaining its other hypotheses? The supplied record does not identify the precise theorem, admissible singular spaces, or intended formula, so this is a statement-only formulation requiring mathematical review.
\subsection{Sources}
\begin{itemize}
\item Supplied local UnsolvedMath registry record, permanent ID \(30002072\), title: \emph{Cheeger--Müller Gluing Without Product Structures}; collection: Oberwolfach Reports Open Problems; exact source problem code: \texttt{catalogue entry}. The record gives proposed year \(2012\) and names no individual proposer.
\item Supplied report citation: \emph{Analysis and Geometric Singularities (2013)}, Oberwolfach Reports, pp.~1487--1562, DOI \texttt{10.4171/owr/2012/25}. \url{https://doi.org/10.4171/owr/2012/25}. The background attributes the problem to workshop open-problem discussions documented in Oberwolfach Reports.
\item Supplied literature source, labelled ``Lesch analytic-torsion paper'': \url{https://doi.org/10.2140/apde.2013.6.221}.
\item Supplied literature source, labelled ``original OWR report'': \url{https://doi.org/10.4171/owr/2012/25}.
\item Literature assessment dated \texttt{2026-08-21}: the supplied curation labels the problem open, states that removal of the product assumption remains unresolved, and states that the published theorem retains that assumption. This is supplied dated curation, not an independently verified literature assessment. The background expressly calls for independent verification of time-sensitive status and cited claims.
\item The registry assigns statement status \texttt{corrected\_verified} and reports that the official report poses exactly the removal of the product-near-boundary hypothesis. This is supplied source-verification metadata; no independent source inspection is asserted here.
\end{itemize}
\end{document}
